% =========================================================================
% Figura (apuntes manuscritos, pág. 11): líneas de campo del cuadrupolo plano
% de cuatro cargas (+q en (+-a,0), -q en (0,+-a)), en el plano z = 0, con la
% paleta del documento. Datos: figuras/datos/cua_plano_*.tex (L = 4a).
% =========================================================================
\begin{figure}[ht]
\centering
\begin{tikzpicture}[x=0.78cm, y=0.78cm]
  \renewcommand{\LC}[1]{\draw[magacero, line width=0.5pt] plot coordinates {#1};}
  \renewcommand{\FL}[3]{\draw[-{Stealth[length=3.6pt,width=3pt]}, magacero, line width=0.5pt] (#1,#2) -- ++(#3:0.01);}
  \fill[magfondo] (-4,-4) rectangle (4,4);
  \begin{scope}\clip (-4,-4) rectangle (4,4);
    \draw[pizarra!60, line width=0.45pt, densely dashed] (-4,-4) -- (4,4) (-4,4) -- (4,-4);
    % Cuadrupolo plano (cuatro cargas) - líneas de campo
% Generado por scripts/generar_datos.py (no editar a mano).
\LC{(1.069,0.012) (1.128,0.023) (1.187,0.033) (1.246,0.044) (1.305,0.055) (1.364,0.066) (1.423,0.078) (1.482,0.090) (1.540,0.102) (1.599,0.114) (1.658,0.127) (1.716,0.141) (1.775,0.155) (1.833,0.169) (1.891,0.184) (1.949,0.199) (2.007,0.215) (2.064,0.232) (2.122,0.249) (2.179,0.266) (2.237,0.284) (2.294,0.303) (2.350,0.323) (2.407,0.342) (2.463,0.363) (2.520,0.384) (2.575,0.406) (2.631,0.428) (2.687,0.451) (2.742,0.475) (2.797,0.499) (2.851,0.524) (2.905,0.549) (2.959,0.576) (3.013,0.603) (3.066,0.630) (3.119,0.658) (3.172,0.687) (3.224,0.717) (3.276,0.747) (3.327,0.778) (3.378,0.810) (3.429,0.842) (3.479,0.875) (3.528,0.909) (3.578,0.943) (3.626,0.978) (3.674,1.014) (3.722,1.051) (3.769,1.088) (3.815,1.126) (3.861,1.165) (3.907,1.204) (3.951,1.244) (3.995,1.285) (4.000,1.292)}
\LC{(1.061,0.035) (1.113,0.065) (1.164,0.096) (1.216,0.126) (1.267,0.158) (1.318,0.190) (1.368,0.223) (1.417,0.256) (1.466,0.291) (1.514,0.327) (1.562,0.364) (1.608,0.402) (1.654,0.441) (1.698,0.482) (1.741,0.523) (1.783,0.566) (1.823,0.611) (1.862,0.657) (1.899,0.704) (1.934,0.753) (1.967,0.802) (1.998,0.854) (2.027,0.906) (2.054,0.960) (2.078,1.015) (2.099,1.071) (2.118,1.128) (2.133,1.186) (2.145,1.245) (2.153,1.304) (2.158,1.364) (2.159,1.424) (2.156,1.484) (2.149,1.544) (2.137,1.602) (2.121,1.660) (2.101,1.717) (2.076,1.771) (2.046,1.823) (2.012,1.873) (1.974,1.919) (1.932,1.962) (1.887,2.001) (1.838,2.036) (1.787,2.067) (1.733,2.094) (1.677,2.116) (1.620,2.133) (1.561,2.146) (1.502,2.154) (1.442,2.159) (1.382,2.159) (1.322,2.155) (1.262,2.148) (1.203,2.137) (1.145,2.122) (1.088,2.105) (1.031,2.085) (0.976,2.061) (0.922,2.035) (0.869,2.007) (0.817,1.977) (0.767,1.944) (0.718,1.909) (0.671,1.873) (0.624,1.835) (0.579,1.795) (0.536,1.753) (0.494,1.711) (0.453,1.667) (0.413,1.622) (0.375,1.576) (0.338,1.529) (0.302,1.481) (0.267,1.432) (0.233,1.382) (0.199,1.332) (0.167,1.282) (0.136,1.231) (0.105,1.179) (0.074,1.128) (0.044,1.076) (0.039,1.067)}
\LC{(1.045,0.054) (1.083,0.100) (1.121,0.146) (1.159,0.193) (1.195,0.241) (1.231,0.289) (1.265,0.338) (1.298,0.389) (1.328,0.440) (1.357,0.493) (1.383,0.547) (1.407,0.602) (1.428,0.658) (1.445,0.716) (1.458,0.774) (1.467,0.834) (1.472,0.893) (1.471,0.953) (1.466,1.013) (1.454,1.072) (1.436,1.129) (1.413,1.184) (1.383,1.236) (1.347,1.285) (1.306,1.328) (1.260,1.367) (1.210,1.399) (1.156,1.426) (1.100,1.446) (1.042,1.461) (0.982,1.469) (0.922,1.472) (0.862,1.470) (0.803,1.463) (0.744,1.452) (0.686,1.436) (0.629,1.417) (0.573,1.395) (0.519,1.370) (0.465,1.342) (0.413,1.313) (0.362,1.281) (0.313,1.247) (0.264,1.213) (0.216,1.176) (0.169,1.139) (0.122,1.102) (0.076,1.064) (0.060,1.051)}
\LC{(1.024,0.066) (1.044,0.122) (1.064,0.179) (1.083,0.236) (1.100,0.293) (1.116,0.351) (1.129,0.410) (1.140,0.469) (1.148,0.528) (1.152,0.588) (1.152,0.648) (1.147,0.708) (1.136,0.767) (1.119,0.824) (1.097,0.880) (1.067,0.932) (1.031,0.980) (0.990,1.023) (0.942,1.060) (0.891,1.091) (0.836,1.115) (0.779,1.133) (0.720,1.145) (0.660,1.151) (0.600,1.152) (0.541,1.149) (0.481,1.142) (0.422,1.132) (0.363,1.119) (0.305,1.103) (0.248,1.086) (0.191,1.068) (0.134,1.048) (0.077,1.028) (0.068,1.025)}
\LC{(1.000,0.070) (1.000,0.130) (0.999,0.190) (0.997,0.250) (0.993,0.310) (0.988,0.370) (0.980,0.429) (0.969,0.488) (0.955,0.546) (0.936,0.603) (0.912,0.658) (0.883,0.711) (0.849,0.760) (0.809,0.805) (0.765,0.845) (0.716,0.880) (0.664,0.910) (0.609,0.934) (0.552,0.953) (0.494,0.968) (0.435,0.979) (0.376,0.987) (0.316,0.993) (0.256,0.997) (0.196,0.999) (0.136,1.000) (0.076,1.000)}
\LC{(0.976,0.066) (0.955,0.122) (0.934,0.178) (0.912,0.234) (0.889,0.289) (0.864,0.344) (0.837,0.398) (0.809,0.451) (0.778,0.502) (0.744,0.551) (0.707,0.599) (0.667,0.643) (0.624,0.685) (0.578,0.724) (0.530,0.759) (0.479,0.792) (0.427,0.822) (0.374,0.849) (0.320,0.875) (0.265,0.899) (0.209,0.922) (0.153,0.944) (0.097,0.965) (0.069,0.975)}
\LC{(0.955,0.054) (0.916,0.099) (0.877,0.145) (0.838,0.190) (0.798,0.235) (0.757,0.279) (0.716,0.323) (0.675,0.366) (0.633,0.409) (0.591,0.452) (0.549,0.495) (0.506,0.537) (0.464,0.580) (0.421,0.622) (0.378,0.664) (0.335,0.705) (0.291,0.746) (0.247,0.787) (0.202,0.827) (0.157,0.867) (0.112,0.906) (0.066,0.945) (0.058,0.951)}
\LC{(0.939,0.035) (0.887,0.065) (0.835,0.095) (0.783,0.124) (0.731,0.153) (0.678,0.183) (0.626,0.213) (0.575,0.243) (0.524,0.276) (0.475,0.310) (0.428,0.348) (0.384,0.389) (0.344,0.433) (0.307,0.480) (0.272,0.529) (0.240,0.580) (0.210,0.631) (0.180,0.683) (0.151,0.736) (0.121,0.788) (0.092,0.840) (0.062,0.892) (0.037,0.936)}
\LC{(0.931,0.012) (0.872,0.023) (0.813,0.033) (0.754,0.043) (0.695,0.053) (0.636,0.064) (0.577,0.075) (0.518,0.086) (0.459,0.099) (0.401,0.115) (0.344,0.134) (0.290,0.160) (0.240,0.193) (0.197,0.235) (0.163,0.284) (0.137,0.338) (0.117,0.394) (0.101,0.452) (0.088,0.511) (0.076,0.570) (0.065,0.629) (0.055,0.688) (0.044,0.747) (0.034,0.806) (0.024,0.865) (0.013,0.924)}
\LC{(0.931,-0.012) (0.872,-0.023) (0.813,-0.033) (0.754,-0.043) (0.695,-0.053) (0.636,-0.064) (0.577,-0.075) (0.518,-0.086) (0.459,-0.099) (0.401,-0.115) (0.344,-0.134) (0.290,-0.160) (0.240,-0.193) (0.197,-0.235) (0.163,-0.284) (0.137,-0.338) (0.117,-0.394) (0.101,-0.452) (0.088,-0.511) (0.076,-0.570) (0.065,-0.629) (0.055,-0.688) (0.044,-0.747) (0.034,-0.806) (0.024,-0.865) (0.013,-0.924)}
\LC{(0.939,-0.035) (0.887,-0.065) (0.835,-0.095) (0.783,-0.124) (0.731,-0.153) (0.678,-0.183) (0.626,-0.213) (0.575,-0.243) (0.524,-0.276) (0.475,-0.310) (0.428,-0.348) (0.384,-0.389) (0.344,-0.433) (0.307,-0.480) (0.272,-0.529) (0.240,-0.580) (0.210,-0.631) (0.180,-0.683) (0.151,-0.736) (0.121,-0.788) (0.092,-0.840) (0.062,-0.892) (0.037,-0.936)}
\LC{(0.955,-0.054) (0.916,-0.099) (0.877,-0.145) (0.838,-0.190) (0.798,-0.235) (0.757,-0.279) (0.716,-0.323) (0.675,-0.366) (0.633,-0.409) (0.591,-0.452) (0.549,-0.495) (0.506,-0.537) (0.464,-0.580) (0.421,-0.622) (0.378,-0.664) (0.335,-0.705) (0.291,-0.746) (0.247,-0.787) (0.202,-0.827) (0.157,-0.867) (0.112,-0.906) (0.066,-0.945) (0.058,-0.951)}
\LC{(0.976,-0.066) (0.955,-0.122) (0.934,-0.178) (0.912,-0.234) (0.889,-0.289) (0.864,-0.344) (0.837,-0.398) (0.809,-0.451) (0.778,-0.502) (0.744,-0.551) (0.707,-0.599) (0.667,-0.643) (0.624,-0.685) (0.578,-0.724) (0.530,-0.759) (0.479,-0.792) (0.427,-0.822) (0.374,-0.849) (0.320,-0.875) (0.265,-0.899) (0.209,-0.922) (0.153,-0.944) (0.097,-0.965) (0.069,-0.975)}
\LC{(1.000,-0.070) (1.000,-0.130) (0.999,-0.190) (0.997,-0.250) (0.993,-0.310) (0.988,-0.370) (0.980,-0.429) (0.969,-0.488) (0.955,-0.546) (0.936,-0.603) (0.912,-0.658) (0.883,-0.711) (0.849,-0.760) (0.809,-0.805) (0.765,-0.845) (0.716,-0.880) (0.664,-0.910) (0.609,-0.934) (0.552,-0.953) (0.494,-0.968) (0.435,-0.979) (0.376,-0.987) (0.316,-0.993) (0.256,-0.997) (0.196,-0.999) (0.136,-1.000) (0.076,-1.000)}
\LC{(1.024,-0.066) (1.044,-0.122) (1.064,-0.179) (1.083,-0.236) (1.100,-0.293) (1.116,-0.351) (1.129,-0.410) (1.140,-0.469) (1.148,-0.528) (1.152,-0.588) (1.152,-0.648) (1.147,-0.708) (1.136,-0.767) (1.119,-0.824) (1.097,-0.880) (1.067,-0.932) (1.031,-0.980) (0.990,-1.023) (0.942,-1.060) (0.891,-1.091) (0.836,-1.115) (0.779,-1.133) (0.720,-1.145) (0.660,-1.151) (0.600,-1.152) (0.541,-1.149) (0.481,-1.142) (0.422,-1.132) (0.363,-1.119) (0.305,-1.103) (0.248,-1.086) (0.191,-1.068) (0.134,-1.048) (0.077,-1.028) (0.068,-1.025)}
\LC{(1.045,-0.054) (1.083,-0.100) (1.121,-0.146) (1.159,-0.193) (1.195,-0.241) (1.231,-0.289) (1.265,-0.338) (1.298,-0.389) (1.328,-0.440) (1.357,-0.493) (1.383,-0.547) (1.407,-0.602) (1.428,-0.658) (1.445,-0.716) (1.458,-0.774) (1.467,-0.834) (1.472,-0.893) (1.471,-0.953) (1.466,-1.013) (1.454,-1.072) (1.436,-1.129) (1.413,-1.184) (1.383,-1.236) (1.347,-1.285) (1.306,-1.328) (1.260,-1.367) (1.210,-1.399) (1.156,-1.426) (1.100,-1.446) (1.042,-1.461) (0.982,-1.469) (0.922,-1.472) (0.862,-1.470) (0.803,-1.463) (0.744,-1.452) (0.686,-1.436) (0.629,-1.417) (0.573,-1.395) (0.519,-1.370) (0.465,-1.342) (0.413,-1.313) (0.362,-1.281) (0.313,-1.247) (0.264,-1.213) (0.216,-1.176) (0.169,-1.139) (0.122,-1.102) (0.076,-1.064) (0.060,-1.051)}
\LC{(1.061,-0.035) (1.113,-0.065) (1.164,-0.096) (1.216,-0.126) (1.267,-0.158) (1.318,-0.190) (1.368,-0.223) (1.417,-0.256) (1.466,-0.291) (1.514,-0.327) (1.562,-0.364) (1.608,-0.402) (1.654,-0.441) (1.698,-0.482) (1.741,-0.523) (1.783,-0.566) (1.823,-0.611) (1.862,-0.657) (1.899,-0.704) (1.934,-0.753) (1.967,-0.802) (1.998,-0.854) (2.027,-0.906) (2.054,-0.960) (2.078,-1.015) (2.099,-1.071) (2.118,-1.128) (2.133,-1.186) (2.145,-1.245) (2.153,-1.304) (2.158,-1.364) (2.159,-1.424) (2.156,-1.484) (2.149,-1.544) (2.137,-1.602) (2.121,-1.660) (2.101,-1.717) (2.076,-1.771) (2.046,-1.823) (2.012,-1.873) (1.974,-1.919) (1.932,-1.962) (1.887,-2.001) (1.838,-2.036) (1.787,-2.067) (1.733,-2.094) (1.677,-2.116) (1.620,-2.133) (1.561,-2.146) (1.502,-2.154) (1.442,-2.159) (1.382,-2.159) (1.322,-2.155) (1.262,-2.148) (1.203,-2.137) (1.145,-2.122) (1.088,-2.105) (1.031,-2.085) (0.976,-2.061) (0.922,-2.035) (0.869,-2.007) (0.817,-1.977) (0.767,-1.944) (0.718,-1.909) (0.671,-1.873) (0.624,-1.835) (0.579,-1.795) (0.536,-1.753) (0.494,-1.711) (0.453,-1.667) (0.413,-1.622) (0.375,-1.576) (0.338,-1.529) (0.302,-1.481) (0.267,-1.432) (0.233,-1.382) (0.199,-1.332) (0.167,-1.282) (0.136,-1.231) (0.105,-1.179) (0.074,-1.128) (0.044,-1.076) (0.039,-1.067)}
\LC{(1.069,-0.012) (1.128,-0.023) (1.187,-0.033) (1.246,-0.044) (1.305,-0.055) (1.364,-0.066) (1.423,-0.078) (1.482,-0.090) (1.540,-0.102) (1.599,-0.114) (1.658,-0.127) (1.716,-0.141) (1.775,-0.155) (1.833,-0.169) (1.891,-0.184) (1.949,-0.199) (2.007,-0.215) (2.064,-0.232) (2.122,-0.249) (2.179,-0.266) (2.237,-0.284) (2.294,-0.303) (2.350,-0.323) (2.407,-0.342) (2.463,-0.363) (2.520,-0.384) (2.575,-0.406) (2.631,-0.428) (2.687,-0.451) (2.742,-0.475) (2.797,-0.499) (2.851,-0.524) (2.905,-0.549) (2.959,-0.576) (3.013,-0.603) (3.066,-0.630) (3.119,-0.658) (3.172,-0.687) (3.224,-0.717) (3.276,-0.747) (3.327,-0.778) (3.378,-0.810) (3.429,-0.842) (3.479,-0.875) (3.528,-0.909) (3.578,-0.943) (3.626,-0.978) (3.674,-1.014) (3.722,-1.051) (3.769,-1.088) (3.815,-1.126) (3.861,-1.165) (3.907,-1.204) (3.951,-1.244) (3.995,-1.285) (4.000,-1.292)}
\LC{(-0.931,0.012) (-0.872,0.023) (-0.813,0.033) (-0.754,0.043) (-0.695,0.053) (-0.636,0.064) (-0.577,0.075) (-0.518,0.086) (-0.459,0.099) (-0.401,0.115) (-0.344,0.134) (-0.290,0.160) (-0.240,0.193) (-0.197,0.235) (-0.163,0.284) (-0.137,0.338) (-0.117,0.394) (-0.101,0.452) (-0.088,0.511) (-0.076,0.570) (-0.065,0.629) (-0.055,0.688) (-0.044,0.747) (-0.034,0.806) (-0.024,0.865) (-0.013,0.924)}
\LC{(-0.939,0.035) (-0.887,0.065) (-0.835,0.095) (-0.783,0.124) (-0.731,0.153) (-0.678,0.183) (-0.626,0.213) (-0.575,0.243) (-0.524,0.276) (-0.475,0.310) (-0.428,0.348) (-0.384,0.389) (-0.344,0.433) (-0.307,0.480) (-0.272,0.529) (-0.240,0.580) (-0.210,0.631) (-0.180,0.683) (-0.151,0.736) (-0.121,0.788) (-0.092,0.840) (-0.062,0.892) (-0.037,0.936)}
\LC{(-0.955,0.054) (-0.916,0.099) (-0.877,0.145) (-0.838,0.190) (-0.798,0.235) (-0.757,0.279) (-0.716,0.323) (-0.675,0.366) (-0.633,0.409) (-0.591,0.452) (-0.549,0.495) (-0.506,0.537) (-0.464,0.580) (-0.421,0.622) (-0.378,0.664) (-0.335,0.705) (-0.291,0.746) (-0.247,0.787) (-0.202,0.827) (-0.157,0.867) (-0.112,0.906) (-0.066,0.945) (-0.058,0.951)}
\LC{(-0.976,0.066) (-0.955,0.122) (-0.934,0.178) (-0.912,0.234) (-0.889,0.289) (-0.864,0.344) (-0.837,0.398) (-0.809,0.451) (-0.778,0.502) (-0.744,0.551) (-0.707,0.599) (-0.667,0.643) (-0.624,0.685) (-0.578,0.724) (-0.530,0.759) (-0.479,0.792) (-0.427,0.822) (-0.374,0.849) (-0.320,0.875) (-0.265,0.899) (-0.209,0.922) (-0.153,0.944) (-0.097,0.965) (-0.069,0.975)}
\LC{(-1.000,0.070) (-1.000,0.130) (-0.999,0.190) (-0.997,0.250) (-0.993,0.310) (-0.988,0.370) (-0.980,0.429) (-0.969,0.488) (-0.955,0.546) (-0.936,0.603) (-0.912,0.658) (-0.883,0.711) (-0.849,0.760) (-0.809,0.805) (-0.765,0.845) (-0.716,0.880) (-0.664,0.910) (-0.609,0.934) (-0.552,0.953) (-0.494,0.968) (-0.435,0.979) (-0.376,0.987) (-0.316,0.993) (-0.256,0.997) (-0.196,0.999) (-0.136,1.000) (-0.076,1.000)}
\LC{(-1.024,0.066) (-1.044,0.122) (-1.064,0.179) (-1.083,0.236) (-1.100,0.293) (-1.116,0.351) (-1.129,0.410) (-1.140,0.469) (-1.148,0.528) (-1.152,0.588) (-1.152,0.648) (-1.147,0.708) (-1.136,0.767) (-1.119,0.824) (-1.097,0.880) (-1.067,0.932) (-1.031,0.980) (-0.990,1.023) (-0.942,1.060) (-0.891,1.091) (-0.836,1.115) (-0.779,1.133) (-0.720,1.145) (-0.660,1.151) (-0.600,1.152) (-0.541,1.149) (-0.481,1.142) (-0.422,1.132) (-0.363,1.119) (-0.305,1.103) (-0.248,1.086) (-0.191,1.068) (-0.134,1.048) (-0.077,1.028) (-0.068,1.025)}
\LC{(-1.045,0.054) (-1.083,0.100) (-1.121,0.146) (-1.159,0.193) (-1.195,0.241) (-1.231,0.289) (-1.265,0.338) (-1.298,0.389) (-1.328,0.440) (-1.357,0.493) (-1.383,0.547) (-1.407,0.602) (-1.428,0.658) (-1.445,0.716) (-1.458,0.774) (-1.467,0.834) (-1.472,0.893) (-1.471,0.953) (-1.466,1.013) (-1.454,1.072) (-1.436,1.129) (-1.413,1.184) (-1.383,1.236) (-1.347,1.285) (-1.306,1.328) (-1.260,1.367) (-1.210,1.399) (-1.156,1.426) (-1.100,1.446) (-1.042,1.461) (-0.982,1.469) (-0.922,1.472) (-0.862,1.470) (-0.803,1.463) (-0.744,1.452) (-0.686,1.436) (-0.629,1.417) (-0.573,1.395) (-0.519,1.370) (-0.465,1.342) (-0.413,1.313) (-0.362,1.281) (-0.313,1.247) (-0.264,1.213) (-0.216,1.176) (-0.169,1.139) (-0.122,1.102) (-0.076,1.064) (-0.060,1.051)}
\LC{(-1.061,0.035) (-1.113,0.065) (-1.164,0.096) (-1.216,0.126) (-1.267,0.158) (-1.318,0.190) (-1.368,0.223) (-1.417,0.256) (-1.466,0.291) (-1.514,0.327) (-1.562,0.364) (-1.608,0.402) (-1.654,0.441) (-1.698,0.482) (-1.741,0.523) (-1.783,0.566) (-1.823,0.611) (-1.862,0.657) (-1.899,0.704) (-1.934,0.753) (-1.967,0.802) (-1.998,0.854) (-2.027,0.906) (-2.054,0.960) (-2.078,1.015) (-2.099,1.071) (-2.118,1.128) (-2.133,1.186) (-2.145,1.245) (-2.153,1.304) (-2.158,1.364) (-2.159,1.424) (-2.156,1.484) (-2.149,1.544) (-2.137,1.602) (-2.121,1.660) (-2.101,1.717) (-2.076,1.771) (-2.046,1.823) (-2.012,1.873) (-1.974,1.919) (-1.932,1.962) (-1.887,2.001) (-1.838,2.036) (-1.787,2.067) (-1.733,2.094) (-1.677,2.116) (-1.620,2.133) (-1.561,2.146) (-1.502,2.154) (-1.442,2.159) (-1.382,2.159) (-1.322,2.155) (-1.262,2.148) (-1.203,2.137) (-1.145,2.122) (-1.088,2.105) (-1.031,2.085) (-0.976,2.061) (-0.922,2.035) (-0.869,2.007) (-0.817,1.977) (-0.767,1.944) (-0.718,1.909) (-0.671,1.873) (-0.624,1.835) (-0.579,1.795) (-0.536,1.753) (-0.494,1.711) (-0.453,1.667) (-0.413,1.622) (-0.375,1.576) (-0.338,1.529) (-0.302,1.481) (-0.267,1.432) (-0.233,1.382) (-0.199,1.332) (-0.167,1.282) (-0.136,1.231) (-0.105,1.179) (-0.074,1.128) (-0.044,1.076) (-0.039,1.067)}
\LC{(-1.069,0.012) (-1.128,0.023) (-1.187,0.033) (-1.246,0.044) (-1.305,0.055) (-1.364,0.066) (-1.423,0.078) (-1.482,0.090) (-1.540,0.102) (-1.599,0.114) (-1.658,0.127) (-1.716,0.141) (-1.775,0.155) (-1.833,0.169) (-1.891,0.184) (-1.949,0.199) (-2.007,0.215) (-2.064,0.232) (-2.122,0.249) (-2.179,0.266) (-2.237,0.284) (-2.294,0.303) (-2.350,0.323) (-2.407,0.342) (-2.463,0.363) (-2.520,0.384) (-2.575,0.406) (-2.631,0.428) (-2.687,0.451) (-2.742,0.475) (-2.797,0.499) (-2.851,0.524) (-2.905,0.549) (-2.959,0.576) (-3.013,0.603) (-3.066,0.630) (-3.119,0.658) (-3.172,0.687) (-3.224,0.717) (-3.276,0.747) (-3.327,0.778) (-3.378,0.810) (-3.429,0.842) (-3.479,0.875) (-3.528,0.909) (-3.578,0.943) (-3.626,0.978) (-3.674,1.014) (-3.722,1.051) (-3.769,1.088) (-3.815,1.126) (-3.861,1.165) (-3.907,1.204) (-3.951,1.244) (-3.995,1.285) (-4.000,1.292)}
\LC{(-1.069,-0.012) (-1.128,-0.023) (-1.187,-0.033) (-1.246,-0.044) (-1.305,-0.055) (-1.364,-0.066) (-1.423,-0.078) (-1.482,-0.090) (-1.540,-0.102) (-1.599,-0.114) (-1.658,-0.127) (-1.716,-0.141) (-1.775,-0.155) (-1.833,-0.169) (-1.891,-0.184) (-1.949,-0.199) (-2.007,-0.215) (-2.064,-0.232) (-2.122,-0.249) (-2.179,-0.266) (-2.237,-0.284) (-2.294,-0.303) (-2.350,-0.323) (-2.407,-0.342) (-2.463,-0.363) (-2.520,-0.384) (-2.575,-0.406) (-2.631,-0.428) (-2.687,-0.451) (-2.742,-0.475) (-2.797,-0.499) (-2.851,-0.524) (-2.905,-0.549) (-2.959,-0.576) (-3.013,-0.603) (-3.066,-0.630) (-3.119,-0.658) (-3.172,-0.687) (-3.224,-0.717) (-3.276,-0.747) (-3.327,-0.778) (-3.378,-0.810) (-3.429,-0.842) (-3.479,-0.875) (-3.528,-0.909) (-3.578,-0.943) (-3.626,-0.978) (-3.674,-1.014) (-3.722,-1.051) (-3.769,-1.088) (-3.815,-1.126) (-3.861,-1.165) (-3.907,-1.204) (-3.951,-1.244) (-3.995,-1.285) (-4.000,-1.292)}
\LC{(-1.061,-0.035) (-1.113,-0.065) (-1.164,-0.096) (-1.216,-0.126) (-1.267,-0.158) (-1.318,-0.190) (-1.368,-0.223) (-1.417,-0.256) (-1.466,-0.291) (-1.514,-0.327) (-1.562,-0.364) (-1.608,-0.402) (-1.654,-0.441) (-1.698,-0.482) (-1.741,-0.523) (-1.783,-0.566) (-1.823,-0.611) (-1.862,-0.657) (-1.899,-0.704) (-1.934,-0.753) (-1.967,-0.802) (-1.998,-0.854) (-2.027,-0.906) (-2.054,-0.960) (-2.078,-1.015) (-2.099,-1.071) (-2.118,-1.128) (-2.133,-1.186) (-2.145,-1.245) (-2.153,-1.304) (-2.158,-1.364) (-2.159,-1.424) (-2.156,-1.484) (-2.149,-1.544) (-2.137,-1.602) (-2.121,-1.660) (-2.101,-1.717) (-2.076,-1.771) (-2.046,-1.823) (-2.012,-1.873) (-1.974,-1.919) (-1.932,-1.962) (-1.887,-2.001) (-1.838,-2.036) (-1.787,-2.067) (-1.733,-2.094) (-1.677,-2.116) (-1.620,-2.133) (-1.561,-2.146) (-1.502,-2.154) (-1.442,-2.159) (-1.382,-2.159) (-1.322,-2.155) (-1.262,-2.148) (-1.203,-2.137) (-1.145,-2.122) (-1.088,-2.105) (-1.031,-2.085) (-0.976,-2.061) (-0.922,-2.035) (-0.869,-2.007) (-0.817,-1.977) (-0.767,-1.944) (-0.718,-1.909) (-0.671,-1.873) (-0.624,-1.835) (-0.579,-1.795) (-0.536,-1.753) (-0.494,-1.711) (-0.453,-1.667) (-0.413,-1.622) (-0.375,-1.576) (-0.338,-1.529) (-0.302,-1.481) (-0.267,-1.432) (-0.233,-1.382) (-0.199,-1.332) (-0.167,-1.282) (-0.136,-1.231) (-0.105,-1.179) (-0.074,-1.128) (-0.044,-1.076) (-0.039,-1.067)}
\LC{(-1.045,-0.054) (-1.083,-0.100) (-1.121,-0.146) (-1.159,-0.193) (-1.195,-0.241) (-1.231,-0.289) (-1.265,-0.338) (-1.298,-0.389) (-1.328,-0.440) (-1.357,-0.493) (-1.383,-0.547) (-1.407,-0.602) (-1.428,-0.658) (-1.445,-0.716) (-1.458,-0.774) (-1.467,-0.834) (-1.472,-0.893) (-1.471,-0.953) (-1.466,-1.013) (-1.454,-1.072) (-1.436,-1.129) (-1.413,-1.184) (-1.383,-1.236) (-1.347,-1.285) (-1.306,-1.328) (-1.260,-1.367) (-1.210,-1.399) (-1.156,-1.426) (-1.100,-1.446) (-1.042,-1.461) (-0.982,-1.469) (-0.922,-1.472) (-0.862,-1.470) (-0.803,-1.463) (-0.744,-1.452) (-0.686,-1.436) (-0.629,-1.417) (-0.573,-1.395) (-0.519,-1.370) (-0.465,-1.342) (-0.413,-1.313) (-0.362,-1.281) (-0.313,-1.247) (-0.264,-1.213) (-0.216,-1.176) (-0.169,-1.139) (-0.122,-1.102) (-0.076,-1.064) (-0.060,-1.051)}
\LC{(-1.024,-0.066) (-1.044,-0.122) (-1.064,-0.179) (-1.083,-0.236) (-1.100,-0.293) (-1.116,-0.351) (-1.129,-0.410) (-1.140,-0.469) (-1.148,-0.528) (-1.152,-0.588) (-1.152,-0.648) (-1.147,-0.708) (-1.136,-0.767) (-1.119,-0.824) (-1.097,-0.880) (-1.067,-0.932) (-1.031,-0.980) (-0.990,-1.023) (-0.942,-1.060) (-0.891,-1.091) (-0.836,-1.115) (-0.779,-1.133) (-0.720,-1.145) (-0.660,-1.151) (-0.600,-1.152) (-0.541,-1.149) (-0.481,-1.142) (-0.422,-1.132) (-0.363,-1.119) (-0.305,-1.103) (-0.248,-1.086) (-0.191,-1.068) (-0.134,-1.048) (-0.077,-1.028) (-0.068,-1.025)}
\LC{(-1.000,-0.070) (-1.000,-0.130) (-0.999,-0.190) (-0.997,-0.250) (-0.993,-0.310) (-0.988,-0.370) (-0.980,-0.429) (-0.969,-0.488) (-0.955,-0.546) (-0.936,-0.603) (-0.912,-0.658) (-0.883,-0.711) (-0.849,-0.760) (-0.809,-0.805) (-0.765,-0.845) (-0.716,-0.880) (-0.664,-0.910) (-0.609,-0.934) (-0.552,-0.953) (-0.494,-0.968) (-0.435,-0.979) (-0.376,-0.987) (-0.316,-0.993) (-0.256,-0.997) (-0.196,-0.999) (-0.136,-1.000) (-0.076,-1.000)}
\LC{(-0.976,-0.066) (-0.955,-0.122) (-0.934,-0.178) (-0.912,-0.234) (-0.889,-0.289) (-0.864,-0.344) (-0.837,-0.398) (-0.809,-0.451) (-0.778,-0.502) (-0.744,-0.551) (-0.707,-0.599) (-0.667,-0.643) (-0.624,-0.685) (-0.578,-0.724) (-0.530,-0.759) (-0.479,-0.792) (-0.427,-0.822) (-0.374,-0.849) (-0.320,-0.875) (-0.265,-0.899) (-0.209,-0.922) (-0.153,-0.944) (-0.097,-0.965) (-0.069,-0.975)}
\LC{(-0.955,-0.054) (-0.916,-0.099) (-0.877,-0.145) (-0.838,-0.190) (-0.798,-0.235) (-0.757,-0.279) (-0.716,-0.323) (-0.675,-0.366) (-0.633,-0.409) (-0.591,-0.452) (-0.549,-0.495) (-0.506,-0.537) (-0.464,-0.580) (-0.421,-0.622) (-0.378,-0.664) (-0.335,-0.705) (-0.291,-0.746) (-0.247,-0.787) (-0.202,-0.827) (-0.157,-0.867) (-0.112,-0.906) (-0.066,-0.945) (-0.058,-0.951)}
\LC{(-0.939,-0.035) (-0.887,-0.065) (-0.835,-0.095) (-0.783,-0.124) (-0.731,-0.153) (-0.678,-0.183) (-0.626,-0.213) (-0.575,-0.243) (-0.524,-0.276) (-0.475,-0.310) (-0.428,-0.348) (-0.384,-0.389) (-0.344,-0.433) (-0.307,-0.480) (-0.272,-0.529) (-0.240,-0.580) (-0.210,-0.631) (-0.180,-0.683) (-0.151,-0.736) (-0.121,-0.788) (-0.092,-0.840) (-0.062,-0.892) (-0.037,-0.936)}
\LC{(-0.931,-0.012) (-0.872,-0.023) (-0.813,-0.033) (-0.754,-0.043) (-0.695,-0.053) (-0.636,-0.064) (-0.577,-0.075) (-0.518,-0.086) (-0.459,-0.099) (-0.401,-0.115) (-0.344,-0.134) (-0.290,-0.160) (-0.240,-0.193) (-0.197,-0.235) (-0.163,-0.284) (-0.137,-0.338) (-0.117,-0.394) (-0.101,-0.452) (-0.088,-0.511) (-0.076,-0.570) (-0.065,-0.629) (-0.055,-0.688) (-0.044,-0.747) (-0.034,-0.806) (-0.024,-0.865) (-0.013,-0.924)}
\LC{(-2.445,-0.677) (-2.492,-0.715) (-2.538,-0.753) (-2.583,-0.793) (-2.627,-0.834) (-2.670,-0.875) (-2.712,-0.918) (-2.754,-0.962) (-2.794,-1.006) (-2.832,-1.052) (-2.870,-1.099) (-2.906,-1.146) (-2.941,-1.195) (-2.975,-1.245) (-3.007,-1.296) (-3.038,-1.347) (-3.067,-1.400) (-3.094,-1.453) (-3.119,-1.507) (-3.143,-1.563) (-3.165,-1.618) (-3.185,-1.675) (-3.203,-1.732) (-3.218,-1.790) (-3.232,-1.849) (-3.243,-1.908) (-3.252,-1.967) (-3.258,-2.027) (-3.262,-2.087) (-3.263,-2.147) (-3.261,-2.207) (-3.257,-2.266) (-3.249,-2.326) (-3.239,-2.385) (-3.226,-2.444) (-3.210,-2.501) (-3.191,-2.558) (-3.169,-2.614) (-3.144,-2.669) (-3.116,-2.722) (-3.085,-2.773) (-3.051,-2.822) (-3.014,-2.870) (-2.975,-2.915) (-2.933,-2.958) (-2.888,-2.998) (-2.842,-3.036) (-2.793,-3.071) (-2.742,-3.103) (-2.690,-3.133) (-2.636,-3.159) (-2.581,-3.182) (-2.524,-3.203) (-2.467,-3.220) (-2.409,-3.234) (-2.350,-3.246) (-2.290,-3.254) (-2.231,-3.260) (-2.171,-3.262) (-2.111,-3.262) (-2.051,-3.260) (-1.991,-3.255) (-1.932,-3.247) (-1.872,-3.237) (-1.814,-3.224) (-1.756,-3.209) (-1.698,-3.192) (-1.641,-3.173) (-1.585,-3.152) (-1.530,-3.129) (-1.475,-3.104) (-1.421,-3.078) (-1.368,-3.050) (-1.316,-3.020) (-1.265,-2.988) (-1.215,-2.955) (-1.166,-2.921) (-1.118,-2.885) (-1.071,-2.848) (-1.025,-2.809) (-0.979,-2.770) (-0.935,-2.729) (-0.892,-2.687) (-0.850,-2.645) (-0.809,-2.601) (-0.769,-2.556) (-0.730,-2.510) (-0.692,-2.464) (-0.656,-2.416) (-0.620,-2.368) (-0.585,-2.320) (-0.551,-2.270) (-0.518,-2.220) (-0.486,-2.169) (-0.455,-2.118) (-0.426,-2.066) (-0.421,-2.057)}
\LC{(2.872,-0.869) (2.916,-0.908) (2.960,-0.949) (3.004,-0.991) (3.046,-1.033) (3.087,-1.077) (3.127,-1.121) (3.167,-1.167) (3.205,-1.213) (3.242,-1.260) (3.278,-1.308) (3.312,-1.357) (3.346,-1.407) (3.378,-1.458) (3.408,-1.510) (3.438,-1.562) (3.465,-1.615) (3.492,-1.669) (3.516,-1.724) (3.539,-1.779) (3.560,-1.835) (3.580,-1.892) (3.597,-1.950) (3.613,-2.007) (3.627,-2.066) (3.638,-2.125) (3.648,-2.184) (3.655,-2.244) (3.660,-2.303) (3.663,-2.363) (3.663,-2.423) (3.661,-2.483) (3.657,-2.543) (3.650,-2.603) (3.640,-2.662) (3.628,-2.721) (3.613,-2.779) (3.596,-2.836) (3.575,-2.893) (3.553,-2.948) (3.527,-3.002) (3.499,-3.055) (3.468,-3.107) (3.434,-3.156) (3.399,-3.204) (3.360,-3.251) (3.320,-3.295) (3.277,-3.337) (3.232,-3.376) (3.185,-3.414) (3.136,-3.449) (3.085,-3.481) (3.033,-3.511) (2.980,-3.538) (2.925,-3.562) (2.869,-3.584) (2.812,-3.603) (2.755,-3.620) (2.696,-3.634) (2.637,-3.645) (2.578,-3.653) (2.518,-3.659) (2.458,-3.663) (2.398,-3.664) (2.338,-3.662) (2.278,-3.658) (2.219,-3.652) (2.159,-3.644) (2.100,-3.634) (2.042,-3.621) (1.983,-3.607) (1.926,-3.590) (1.869,-3.572) (1.812,-3.552) (1.756,-3.530) (1.701,-3.506) (1.647,-3.481) (1.593,-3.454) (1.540,-3.426) (1.488,-3.396) (1.437,-3.365) (1.386,-3.332) (1.337,-3.298) (1.288,-3.263) (1.241,-3.227) (1.194,-3.189) (1.148,-3.151) (1.103,-3.111) (1.059,-3.070) (1.016,-3.028) (0.973,-2.986) (0.932,-2.942) (0.892,-2.898) (0.852,-2.853) (0.814,-2.807) (0.776,-2.760) (0.740,-2.712) (0.704,-2.664) (0.669,-2.615) (0.635,-2.566) (0.602,-2.516) (0.570,-2.465) (0.539,-2.414) (0.509,-2.362) (0.480,-2.309) (0.451,-2.257) (0.423,-2.203) (0.397,-2.150) (0.371,-2.096) (0.345,-2.041) (0.321,-1.986) (0.305,-1.950)}
\LC{(0.641,-4.000) (0.620,-3.952) (0.598,-3.896) (0.577,-3.839) (0.557,-3.783) (0.537,-3.726) (0.517,-3.670) (0.498,-3.613) (0.479,-3.556) (0.461,-3.499) (0.443,-3.442) (0.425,-3.384) (0.408,-3.327) (0.391,-3.269) (0.375,-3.211) (0.359,-3.154) (0.343,-3.096) (0.328,-3.038) (0.313,-2.980) (0.299,-2.921) (0.284,-2.863) (0.271,-2.805) (0.257,-2.746) (0.244,-2.688) (0.232,-2.629) (0.219,-2.570) (0.207,-2.511) (0.196,-2.453) (0.186,-2.403)}
\LC{(-3.072,0.978) (-3.116,1.019) (-3.158,1.062) (-3.200,1.105) (-3.241,1.149) (-3.281,1.194) (-3.319,1.240) (-3.357,1.286) (-3.393,1.334) (-3.429,1.382) (-3.463,1.432) (-3.496,1.482) (-3.528,1.533) (-3.558,1.584) (-3.587,1.637) (-3.615,1.690) (-3.641,1.744) (-3.665,1.799) (-3.688,1.855) (-3.709,1.911) (-3.729,1.967) (-3.746,2.025) (-3.762,2.083) (-3.776,2.141) (-3.789,2.200) (-3.799,2.259) (-3.807,2.318) (-3.812,2.378) (-3.816,2.438) (-3.818,2.498) (-3.817,2.558) (-3.813,2.618) (-3.808,2.677) (-3.800,2.737) (-3.789,2.796) (-3.776,2.855) (-3.760,2.912) (-3.742,2.970) (-3.721,3.026) (-3.698,3.081) (-3.672,3.135) (-3.643,3.188) (-3.612,3.239) (-3.579,3.289) (-3.543,3.337) (-3.505,3.383) (-3.464,3.428) (-3.422,3.470) (-3.377,3.510) (-3.330,3.548) (-3.282,3.584) (-3.232,3.617) (-3.181,3.647) (-3.128,3.675) (-3.073,3.701) (-3.018,3.724) (-2.961,3.745) (-2.904,3.763) (-2.846,3.778) (-2.788,3.791) (-2.728,3.801) (-2.669,3.809) (-2.609,3.814) (-2.549,3.817) (-2.489,3.817) (-2.429,3.816) (-2.369,3.812) (-2.310,3.806) (-2.250,3.797) (-2.191,3.787) (-2.133,3.775) (-2.074,3.760) (-2.017,3.744) (-1.959,3.726) (-1.903,3.706) (-1.847,3.685) (-1.791,3.662) (-1.737,3.637) (-1.683,3.611) (-1.629,3.583) (-1.577,3.554) (-1.525,3.523) (-1.475,3.491) (-1.425,3.458) (-1.375,3.424) (-1.327,3.388) (-1.280,3.352) (-1.233,3.314) (-1.187,3.275) (-1.142,3.235) (-1.099,3.194) (-1.055,3.152) (-1.013,3.110) (-0.972,3.066) (-0.932,3.022) (-0.892,2.977) (-0.854,2.931) (-0.816,2.884) (-0.779,2.836) (-0.743,2.788) (-0.709,2.740) (-0.674,2.690) (-0.641,2.640) (-0.609,2.590) (-0.577,2.539) (-0.547,2.487) (-0.517,2.435) (-0.488,2.382) (-0.460,2.329) (-0.433,2.276) (-0.407,2.222) (-0.381,2.168) (-0.356,2.113) (-0.332,2.058) (-0.309,2.003) (-0.286,1.947) (-0.275,1.919)}
\LC{(2.400,0.661) (2.447,0.699) (2.493,0.737) (2.538,0.776) (2.582,0.817) (2.626,0.859) (2.668,0.901) (2.709,0.945) (2.749,0.989) (2.788,1.035) (2.825,1.082) (2.862,1.130) (2.897,1.179) (2.930,1.228) (2.962,1.279) (2.992,1.331) (3.021,1.383) (3.048,1.437) (3.074,1.491) (3.097,1.547) (3.119,1.603) (3.138,1.659) (3.155,1.717) (3.171,1.775) (3.184,1.833) (3.194,1.893) (3.202,1.952) (3.208,2.012) (3.211,2.072) (3.211,2.132) (3.208,2.192) (3.203,2.251) (3.195,2.311) (3.183,2.370) (3.169,2.428) (3.152,2.485) (3.131,2.542) (3.108,2.597) (3.081,2.651) (3.052,2.703) (3.019,2.754) (2.984,2.802) (2.946,2.848) (2.905,2.892) (2.862,2.934) (2.816,2.973) (2.768,3.009) (2.719,3.042) (2.667,3.073) (2.613,3.100) (2.559,3.125) (2.503,3.146) (2.445,3.164) (2.387,3.179) (2.329,3.192) (2.269,3.201) (2.210,3.207) (2.150,3.210) (2.090,3.211) (2.030,3.209) (1.970,3.204) (1.911,3.197) (1.851,3.187) (1.793,3.175) (1.734,3.160) (1.677,3.144) (1.620,3.125) (1.564,3.104) (1.508,3.081) (1.453,3.056) (1.400,3.030) (1.347,3.001) (1.295,2.971) (1.244,2.940) (1.194,2.907) (1.144,2.872) (1.096,2.837) (1.049,2.799) (1.003,2.761) (0.958,2.721) (0.914,2.680) (0.871,2.639) (0.829,2.596) (0.789,2.552) (0.749,2.507) (0.710,2.461) (0.672,2.414) (0.636,2.367) (0.600,2.318) (0.565,2.269) (0.532,2.220) (0.499,2.169) (0.468,2.118) (0.442,2.075)}
\LC{(-3.089,1.434) (-3.115,1.488) (-3.139,1.543) (-3.162,1.598) (-3.183,1.654) (-3.201,1.711) (-3.218,1.769) (-3.232,1.827) (-3.244,1.886) (-3.254,1.945) (-3.261,2.005) (-3.266,2.065) (-3.269,2.125) (-3.268,2.185) (-3.265,2.245) (-3.259,2.304) (-3.250,2.364) (-3.238,2.422) (-3.223,2.481) (-3.206,2.538) (-3.185,2.594) (-3.161,2.649) (-3.134,2.703) (-3.105,2.755) (-3.072,2.805) (-3.036,2.854) (-2.998,2.900) (-2.958,2.944) (-2.914,2.986) (-2.869,3.025) (-2.821,3.061) (-2.771,3.094) (-2.720,3.125) (-2.666,3.153) (-2.612,3.178) (-2.556,3.199) (-2.499,3.218) (-2.441,3.234) (-2.382,3.247) (-2.323,3.256) (-2.264,3.263) (-2.204,3.267) (-2.144,3.269) (-2.084,3.267) (-2.024,3.263) (-1.964,3.257) (-1.905,3.248) (-1.846,3.236) (-1.788,3.223) (-1.730,3.207) (-1.673,3.189) (-1.616,3.169) (-1.560,3.147) (-1.505,3.123) (-1.451,3.097) (-1.433,3.088)}
\LC{(-1.162,-4.000) (-1.124,-3.959) (-1.088,-3.912) (-1.052,-3.864) (-1.016,-3.816) (-0.981,-3.767) (-0.947,-3.717) (-0.914,-3.667) (-0.881,-3.617) (-0.849,-3.566) (-0.818,-3.515) (-0.787,-3.464) (-0.757,-3.412) (-0.727,-3.360) (-0.698,-3.307) (-0.670,-3.254) (-0.643,-3.201) (-0.616,-3.147) (-0.589,-3.093) (-0.564,-3.039) (-0.539,-2.985) (-0.514,-2.930) (-0.490,-2.875) (-0.467,-2.819) (-0.445,-2.764) (-0.423,-2.708) (-0.401,-2.652) (-0.380,-2.596) (-0.360,-2.539) (-0.350,-2.511)}
\LC{(-0.859,4.000) (-0.831,3.952) (-0.803,3.899) (-0.775,3.846) (-0.748,3.792) (-0.722,3.738) (-0.696,3.684) (-0.670,3.630) (-0.646,3.575) (-0.621,3.520) (-0.597,3.465) (-0.574,3.410) (-0.551,3.354) (-0.529,3.299) (-0.507,3.243) (-0.486,3.187) (-0.465,3.130) (-0.445,3.074) (-0.425,3.017) (-0.406,2.961) (-0.387,2.904) (-0.369,2.846) (-0.351,2.789) (-0.333,2.732) (-0.325,2.703)}
\LC{(3.348,1.368) (3.382,1.418) (3.414,1.468) (3.445,1.520) (3.474,1.572) (3.502,1.625) (3.529,1.679) (3.554,1.734) (3.577,1.789) (3.598,1.845) (3.618,1.901) (3.636,1.959) (3.652,2.016) (3.667,2.075) (3.679,2.133) (3.689,2.193) (3.697,2.252) (3.703,2.312) (3.706,2.372) (3.708,2.432) (3.706,2.492) (3.703,2.551) (3.697,2.611) (3.688,2.671) (3.677,2.729) (3.663,2.788) (3.647,2.846) (3.628,2.902) (3.606,2.958) (3.582,3.013) (3.555,3.067) (3.525,3.119) (3.493,3.170) (3.458,3.219) (3.421,3.266) (3.382,3.311) (3.340,3.354) (3.296,3.395) (3.250,3.434) (3.202,3.470) (3.153,3.504) (3.102,3.535) (3.049,3.564) (2.995,3.590) (2.940,3.614) (2.884,3.634) (2.826,3.653) (2.768,3.668) (2.710,3.681) (2.651,3.691) (2.591,3.699) (2.532,3.704) (2.472,3.707) (2.412,3.707) (2.352,3.705) (2.292,3.701) (2.232,3.695) (2.173,3.686) (2.114,3.675) (2.055,3.662) (1.997,3.647) (1.939,3.631) (1.882,3.612) (1.826,3.591) (1.770,3.569) (1.715,3.545) (1.661,3.520) (1.607,3.493) (1.555,3.465) (1.503,3.435) (1.451,3.403) (1.401,3.371) (1.368,3.348)}
\LC{(-2.404,-0.022) (-2.464,-0.023) (-2.524,-0.024) (-2.584,-0.026) (-2.644,-0.027) (-2.704,-0.029) (-2.764,-0.030) (-2.824,-0.032) (-2.883,-0.033) (-2.943,-0.035) (-3.003,-0.037) (-3.063,-0.038) (-3.123,-0.040) (-3.183,-0.042) (-3.243,-0.044) (-3.303,-0.046) (-3.363,-0.048) (-3.423,-0.050) (-3.483,-0.052) (-3.543,-0.054) (-3.603,-0.056) (-3.663,-0.058) (-3.723,-0.061) (-3.783,-0.063) (-3.843,-0.065) (-3.903,-0.068) (-3.963,-0.070) (-4.000,-0.072)}
\LC{(-3.229,-1.082) (-3.270,-1.125) (-3.311,-1.168) (-3.351,-1.213) (-3.390,-1.259) (-3.428,-1.305) (-3.465,-1.353) (-3.501,-1.401) (-3.535,-1.450) (-3.569,-1.500) (-3.601,-1.550) (-3.632,-1.602) (-3.661,-1.654) (-3.689,-1.707) (-3.716,-1.761) (-3.741,-1.815) (-3.765,-1.870) (-3.787,-1.926) (-3.808,-1.983) (-3.826,-2.040) (-3.843,-2.097) (-3.859,-2.155) (-3.872,-2.214) (-3.883,-2.273) (-3.893,-2.332) (-3.900,-2.391) (-3.905,-2.451) (-3.908,-2.511) (-3.909,-2.571) (-3.908,-2.631) (-3.904,-2.691) (-3.898,-2.751) (-3.890,-2.810) (-3.879,-2.869) (-3.866,-2.927) (-3.850,-2.985) (-3.831,-3.042) (-3.811,-3.099) (-3.787,-3.154) (-3.761,-3.208) (-3.733,-3.261) (-3.702,-3.312) (-3.669,-3.362) (-3.633,-3.411) (-3.595,-3.457) (-3.555,-3.502) (-3.513,-3.545) (-3.469,-3.585) (-3.423,-3.624) (-3.375,-3.660) (-3.326,-3.694) (-3.275,-3.725) (-3.222,-3.754) (-3.168,-3.781) (-3.113,-3.805) (-3.057,-3.826) (-3.000,-3.845) (-2.943,-3.862) (-2.884,-3.876) (-2.825,-3.887) (-2.766,-3.896) (-2.707,-3.903) (-2.647,-3.907) (-2.587,-3.909) (-2.527,-3.909) (-2.467,-3.906) (-2.407,-3.902) (-2.347,-3.895) (-2.288,-3.886) (-2.229,-3.875) (-2.170,-3.862) (-2.112,-3.848) (-2.055,-3.831) (-1.997,-3.813) (-1.941,-3.793) (-1.885,-3.771) (-1.830,-3.748) (-1.775,-3.723) (-1.721,-3.697) (-1.668,-3.669) (-1.615,-3.640) (-1.564,-3.609) (-1.513,-3.577) (-1.463,-3.544) (-1.414,-3.510) (-1.365,-3.474) (-1.318,-3.438) (-1.271,-3.400) (-1.225,-3.362) (-1.180,-3.322) (-1.136,-3.281) (-1.093,-3.240) (-1.078,-3.226)}
\LC{(1.209,4.000) (1.171,3.959) (1.133,3.913) (1.095,3.866) (1.059,3.818) (1.023,3.770) (0.987,3.722) (0.953,3.673) (0.919,3.623) (0.886,3.573) (0.853,3.523) (0.821,3.472) (0.790,3.421) (0.759,3.369) (0.729,3.317) (0.700,3.265) (0.671,3.212) (0.643,3.159) (0.616,3.106) (0.589,3.052) (0.563,2.998) (0.538,2.944) (0.513,2.889) (0.489,2.834) (0.466,2.779) (0.443,2.723) (0.421,2.668) (0.399,2.612) (0.378,2.556) (0.357,2.499)}
\LC{(2.398,-0.019) (2.458,-0.020) (2.518,-0.022) (2.578,-0.023) (2.638,-0.024) (2.698,-0.025) (2.758,-0.027) (2.818,-0.028) (2.878,-0.030) (2.937,-0.031) (2.997,-0.033) (3.057,-0.034) (3.117,-0.036) (3.177,-0.038) (3.237,-0.039) (3.297,-0.041) (3.357,-0.043) (3.417,-0.044) (3.477,-0.046) (3.537,-0.048) (3.597,-0.050) (3.657,-0.052) (3.717,-0.054) (3.777,-0.056) (3.837,-0.058) (3.897,-0.060) (3.957,-0.063) (4.000,-0.064)}
\LC{(-2.508,-1.142) (-2.535,-1.195) (-2.560,-1.250) (-2.583,-1.305) (-2.603,-1.362) (-2.621,-1.419) (-2.636,-1.477) (-2.648,-1.536) (-2.658,-1.595) (-2.664,-1.655) (-2.668,-1.715) (-2.668,-1.775) (-2.665,-1.835) (-2.658,-1.894) (-2.648,-1.953) (-2.634,-2.012) (-2.617,-2.069) (-2.596,-2.125) (-2.571,-2.180) (-2.543,-2.233) (-2.511,-2.284) (-2.476,-2.332) (-2.437,-2.378) (-2.395,-2.421) (-2.350,-2.461) (-2.303,-2.498) (-2.253,-2.531) (-2.201,-2.561) (-2.147,-2.587) (-2.091,-2.609) (-2.034,-2.628) (-1.976,-2.643) (-1.917,-2.655) (-1.858,-2.663) (-1.798,-2.667) (-1.738,-2.668) (-1.678,-2.666) (-1.618,-2.661) (-1.559,-2.652) (-1.500,-2.641) (-1.442,-2.627) (-1.384,-2.610) (-1.327,-2.591) (-1.271,-2.569) (-1.216,-2.545) (-1.162,-2.519) (-1.136,-2.505)}
\LC{(2.686,-1.108) (2.718,-1.158) (2.749,-1.209) (2.778,-1.262) (2.806,-1.315) (2.831,-1.370) (2.855,-1.425) (2.876,-1.481) (2.895,-1.538) (2.912,-1.595) (2.926,-1.654) (2.938,-1.712) (2.948,-1.772) (2.954,-1.831) (2.958,-1.891) (2.959,-1.951) (2.957,-2.011) (2.952,-2.071) (2.944,-2.130) (2.933,-2.189) (2.918,-2.247) (2.900,-2.304) (2.879,-2.361) (2.854,-2.415) (2.826,-2.468) (2.795,-2.520) (2.760,-2.569) (2.723,-2.616) (2.683,-2.660) (2.640,-2.702) (2.594,-2.741) (2.546,-2.777) (2.496,-2.810) (2.444,-2.839) (2.390,-2.866) (2.335,-2.889) (2.278,-2.909) (2.220,-2.925) (2.162,-2.938) (2.102,-2.948) (2.043,-2.955) (1.983,-2.959) (1.923,-2.959) (1.863,-2.957) (1.803,-2.952) (1.744,-2.944) (1.685,-2.933) (1.626,-2.920) (1.568,-2.904) (1.511,-2.886) (1.455,-2.866) (1.399,-2.844) (1.344,-2.819) (1.290,-2.793) (1.237,-2.765) (1.185,-2.735) (1.134,-2.703) (1.101,-2.681)}
\LC{(2.560,1.233) (2.583,1.288) (2.604,1.344) (2.623,1.401) (2.639,1.459) (2.653,1.517) (2.664,1.576) (2.671,1.636) (2.676,1.695) (2.678,1.755) (2.676,1.815) (2.671,1.875) (2.662,1.934) (2.650,1.993) (2.635,2.051) (2.615,2.108) (2.592,2.163) (2.566,2.217) (2.535,2.269) (2.502,2.318) (2.464,2.366) (2.424,2.410) (2.381,2.451) (2.334,2.490) (2.286,2.524) (2.235,2.556) (2.181,2.584) (2.127,2.608) (2.070,2.629) (2.013,2.646) (1.954,2.659) (1.895,2.669) (1.835,2.675) (1.775,2.678) (1.715,2.677) (1.655,2.673) (1.596,2.666) (1.537,2.657) (1.478,2.644) (1.420,2.629) (1.363,2.611) (1.306,2.591) (1.251,2.568) (1.233,2.560)}
\LC{(0.253,4.000) (0.244,3.947) (0.236,3.888) (0.227,3.828) (0.219,3.769) (0.211,3.710) (0.203,3.650) (0.195,3.591) (0.188,3.531) (0.181,3.471) (0.173,3.412) (0.166,3.352) (0.159,3.293) (0.153,3.233) (0.146,3.173) (0.140,3.114) (0.134,3.054) (0.127,2.994) (0.122,2.935) (0.116,2.875) (0.110,2.815) (0.105,2.755) (0.099,2.696) (0.097,2.666)}
\LC{(-2.303,0.784) (-2.344,0.829) (-2.383,0.875) (-2.420,0.921) (-2.456,0.969) (-2.491,1.019) (-2.523,1.069) (-2.554,1.120) (-2.583,1.173) (-2.611,1.226) (-2.636,1.281) (-2.659,1.336) (-2.679,1.392) (-2.698,1.449) (-2.714,1.507) (-2.727,1.566) (-2.737,1.625) (-2.745,1.684) (-2.749,1.744) (-2.751,1.804) (-2.749,1.864) (-2.744,1.924) (-2.736,1.983) (-2.724,2.042) (-2.709,2.100) (-2.690,2.157) (-2.668,2.213) (-2.642,2.267) (-2.612,2.319) (-2.579,2.369) (-2.543,2.417) (-2.504,2.462) (-2.461,2.505) (-2.416,2.544) (-2.368,2.580) (-2.318,2.613) (-2.266,2.642) (-2.211,2.668) (-2.156,2.691) (-2.099,2.709) (-2.041,2.725) (-1.982,2.736) (-1.922,2.745) (-1.863,2.749) (-1.803,2.751) (-1.743,2.749) (-1.683,2.745) (-1.623,2.737) (-1.564,2.726) (-1.506,2.713) (-1.448,2.697) (-1.391,2.679) (-1.335,2.658) (-1.279,2.635) (-1.225,2.610) (-1.171,2.583) (-1.119,2.554) (-1.068,2.522) (-1.017,2.490) (-0.968,2.455) (-0.920,2.419) (-0.873,2.382) (-0.828,2.343) (-0.798,2.316)}

    % Cuadrupolo plano (cuatro cargas) - puntas de flecha
% Generado por scripts/generar_datos.py (no editar a mano).
\FL{2.640}{0.432}{22.3}
\FL{1.947}{1.948}{135.1}
\FL{1.321}{1.314}{134.5}
\FL{1.004}{1.009}{135.5}
\FL{0.809}{0.805}{134.7}
\FL{0.653}{0.658}{135.2}
\FL{0.520}{0.523}{135.0}
\FL{0.384}{0.389}{134.8}
\FL{0.218}{0.213}{135.7}
\FL{0.218}{-0.213}{-135.7}
\FL{0.384}{-0.389}{-134.8}
\FL{0.520}{-0.523}{-135.0}
\FL{0.653}{-0.658}{-135.2}
\FL{0.809}{-0.805}{-134.7}
\FL{1.004}{-1.009}{-135.5}
\FL{1.321}{-1.314}{-134.5}
\FL{1.947}{-1.948}{-135.1}
\FL{2.640}{-0.432}{-22.3}
\FL{-0.218}{0.213}{44.3}
\FL{-0.384}{0.389}{45.2}
\FL{-0.520}{0.523}{45.0}
\FL{-0.653}{0.658}{44.8}
\FL{-0.809}{0.805}{45.3}
\FL{-1.004}{1.009}{44.5}
\FL{-1.321}{1.314}{45.5}
\FL{-1.947}{1.948}{44.9}
\FL{-2.640}{0.432}{157.7}
\FL{-2.640}{-0.432}{-157.7}
\FL{-1.947}{-1.948}{-44.9}
\FL{-1.321}{-1.314}{-45.5}
\FL{-1.004}{-1.009}{-44.5}
\FL{-0.809}{-0.805}{-45.3}
\FL{-0.653}{-0.658}{-44.8}
\FL{-0.520}{-0.523}{-45.0}
\FL{-0.384}{-0.389}{-45.2}
\FL{-0.218}{-0.213}{-44.3}
\FL{-3.230}{-1.839}{-102.2}
\FL{-1.395}{-3.064}{28.2}
\FL{3.659}{-2.283}{-85.5}
\FL{1.280}{-3.257}{143.2}
\FL{0.375}{-3.211}{105.7}
\FL{-3.818}{2.498}{90.3}
\FL{-1.248}{3.327}{-39.3}
\FL{3.175}{1.794}{77.2}
\FL{1.417}{3.039}{-153.4}
\FL{-2.943}{2.958}{44.4}
\FL{-0.689}{-3.289}{61.8}
\FL{-0.555}{3.364}{-67.9}
\FL{3.340}{3.354}{135.5}
\FL{-3.203}{-0.043}{-178.2}
\FL{-3.882}{-2.263}{-100.3}
\FL{-2.249}{-3.879}{10.8}
\FL{0.710}{3.283}{-119.1}
\FL{3.207}{-0.038}{-1.6}
\FL{-2.402}{-2.414}{-44.4}
\FL{2.662}{-2.681}{-135.9}
\FL{2.417}{2.417}{135.0}
\FL{0.164}{3.332}{-96.6}
\FL{-2.490}{2.477}{45.6}

  \end{scope}
  \node[cargapos] at (1,0) {$+$};
  \node[cargapos] at (-1,0) {$+$};
  \node[carganeg] at (0,1) {$-$};
  \node[carganeg] at (0,-1) {$-$};
  \colorlet{marcoc}{grispizarra!70}
  \marcoMMA{4}{2}{$x/a$}{$y/a$}
\end{tikzpicture}
\caption{Cuadrupolo plano del \cref{ej:cuadrupolo-plano}: $+q$ en $(\pm a,0,0)$ y $-q$ en $(0,\pm a,0)$, en el plano $z=0$. Líneas de campo (magenta, con su sentido), trazadas desde las cargas positivas a
intervalos angulares iguales y completadas con las líneas que llegan desde lejos; las diagonales $y=\pm x$ (a trazos) forman la equipotencial nula. El campo se «estira» según $x$ (donde están las cargas positivas, $Q_{xx}>0$) y se «comprime» según
$y$ ($Q_{yy}<0$).}
\label{fig:cuadrupolo-plano}
\end{figure}
